% META: {"id": "1NSI-ADGK-ME-001", "chapitre": "1NSI-ALGO-DICHO-GLOUTON-KNN", "type_objet": "methode", "capacites": ["1NSI-ALGO-DICHO-GLOUTON-KNN-C1"], "capacites_codes": ["C1"], "methodes": ["M1"], "mode_creation": "ex_nihilo", "sources_inspiration": [], "status": "needs_review"}
\begin{fichemethode}{M1}{Écrire une recherche dichotomique et prouver sa terminaison}
\textbf{Quand l'utiliser ?} Quand tu cherches une valeur dans un \textbf{tableau trié} et que tu veux un algorithme bien plus rapide que le parcours séquentiel — ou quand l'énoncé demande de \textbf{prouver la terminaison} avec un variant de boucle.

\textbf{Pas à pas :}
\begin{enumerate}
  \item Vérifie que le tableau est \textbf{trié} : la dichotomie ne fonctionne que sur un tableau trié.
  \item Encadre la zone de recherche par deux indices : \lstinline{g = 0} (gauche) et \lstinline{d = len(t) - 1} (droite).
  \item Tant que \lstinline{g <= d} : calcule le milieu \lstinline{m = (g + d) // 2}.
  \item Compare \lstinline{t[m]} à la valeur cherchée \lstinline{v} :
    si \lstinline{t[m] == v}, c'est trouvé ; si \lstinline{t[m] < v}, la valeur est à droite : \lstinline{g = m + 1} ;
    sinon elle est à gauche : \lstinline{d = m - 1}.
  \item Si la boucle se termine sans trouver, renvoie \lstinline{-1}.
  \item \textbf{Terminaison} : pose le variant $V = d - g$. C'est un entier, il est
    positif ou nul tant que la boucle tourne, et chaque tour le fait
    \textbf{strictement décroître} (le milieu est exclu de la nouvelle zone).
    Une suite d'entiers positifs strictement décroissante est finie : la boucle s'arrête.
\end{enumerate}

\textbf{Exemple :}
\begin{python}
def recherche_dichotomique(t, v):
    g, d = 0, len(t) - 1
    while g <= d:            # variant : d - g
        m = (g + d) // 2
        if t[m] == v:
            return m
        if t[m] < v:
            g = m + 1        # le milieu m est exclu
        else:
            d = m - 1        # le milieu m est exclu
    return -1

t = [2, 5, 7, 11, 13, 17, 19, 23]
recherche_dichotomique(t, 13)   # renvoie 4
recherche_dichotomique(t, 4)    # renvoie -1
\end{python}
Sur \lstinline{t} (8 valeurs), chercher 13 fait décroître le variant $d-g$ :
$7 \to 3 \to 0$ — trois tours suffisent, contre jusqu'à 8 pour un parcours séquentiel.

\textbf{Pièges :}
\begin{itemize}
  \item Appliquer la dichotomie à un tableau \textbf{non trié} : le résultat est faux sans message d'erreur.
  \item Écrire \lstinline{g = m} ou \lstinline{d = m} au lieu de \lstinline{m + 1} / \lstinline{m - 1} : le variant ne décroît plus toujours et la boucle peut devenir infinie.
  \item Confondre \textbf{variant} (entier qui décroît, prouve la terminaison) et \textbf{invariant} (propriété qui reste vraie, prouve la correction).
  \item Oublier le cas du tableau vide : avec \lstinline{g = 0} et \lstinline{d = -1}, la boucle ne s'exécute pas et on renvoie bien \lstinline{-1}.
\end{itemize}

\textbf{Vérifier :} teste sur un tableau vide, un singleton, la première et la dernière case, et une valeur absente ; affiche $d-g$ à chaque tour et contrôle qu'il décroît strictement en restant $\geq 0$.

\textbf{S'entraîner :} exercices C1 du chapitre (dichotomie).
\end{fichemethode}
